% carried verbatim from the v2 tree (hardware.tex); NOT part of this session's deliverable.
% CENTRING, measured 2026-09-19 on the composed page at 600 dpi.  The tikz bounding
% box of this two-panel fragment is not symmetric about its ink: the ink is 421.6 pt
% wide but the box TeX centres is 404.2 pt, because panel (b)'s title overhangs the
% box on the right while panel (a)'s reserved label space on the left is not filled.
% \centering therefore put the graphic 17.3 pt right of the text-block centre, while
% the other sixteen figures sit between -2.9 and +1.3 pt.  \leftskip = -2d plus 1fil
% biases a centred line by exactly -d (here d = 17.3 pt) and changes nothing else;
% measured again after the fix: +0.01 pt.  A \hspace* pair does NOT work here: it
% starts the paragraph before the fragment is read, which turns the end-of-line
% spaces of the fragment's three \definecolor lines into real interword glue and
% delivers only 12.7 pt of the 17.3 pt asked for.  The fragment itself is untouched
% (it is one of the eight byte-identical artefacts check_figures.py regenerates).
\begin{figure*}[tp]\centering
\leftskip=-34.6pt plus 1fil\relax
% fragment: \input into the figure environment -> \normalsize equals the body font exactly.
% AUTO-GENERATED by src/make_hardware_hero.py from data/hw_lucj_n2_result.json -- no hand-typed numbers.
\definecolor{inkPrim}{HTML}{1B1B1B}\definecolor{inkMute}{HTML}{6E6E6E}
\definecolor{warmFill}{HTML}{E8770C}\definecolor{warmLine}{HTML}{7A1E10}
\definecolor{hwWarm}{HTML}{C4360C}\definecolor{simSlate}{HTML}{9AA1AC}\definecolor{band}{HTML}{ECE7DF}
\begin{tikzpicture}[font=\rmfamily]
\pgfplotsset{hpanel/.style={width=7.45cm,height=6.0cm,tick align=outside,tick pos=left,axis line style={inkPrim,line width=0.6pt},
  tick style={inkPrim,line width=0.6pt},tick label style={font=\footnotesize},label style={font=\small},
  title style={font=\small,yshift=1pt},ymajorgrids,major grid style={inkMute,opacity=0.10,line width=0.3pt}}}
% ===================== (a) ibm_fez : spectral function A(omega) =====================
% CLEAN: curve + title only; all technical detail (shots, mitigation, coverage) lives once, in the caption.
\begin{axis}[name=axa,hpanel,
  title={(a)\ \ \texttt{ibm\_fez} $\cdot$ 12 qubits: \ $A(\omega)$},
  xlabel={frequency \ $\omega-E_0$}, ylabel={$A(\omega)$},
  xmin=2.31,xmax=21.02,ymin=0,ymax=0.603]
  \addplot[draw=none,fill=warmFill,fill opacity=0.28] table[x=w,y=A]{figs/heron_hot.dat} \closedcycle;
  \addplot[warmLine,line width=1.3pt] table[x=w,y=A]{figs/heron_hot.dat};
\end{axis}
% ===================== (b) ibm_marrakesh : N2 ground-state energy =====================
% CLEAN: two curves + which-is-which + the two headline numbers; specs/jobs live once, in the caption.
\begin{axis}[name=axb,hpanel,at={(axa.south east)},anchor=south west,xshift=1.4cm,
  title={(b)\ \ \texttt{ibm\_marrakesh} $\cdot$ 24 qubits: \ $\mathrm{N_2}$ energy},
  xlabel={configuration-recovery step}, ylabel={error to FCI \ (mHa)},
  ymode=log,xmin=-0.35,xmax=4.55,ymin=0.28,ymax=90,
  xtick={0,1,2,3,4},ytick={0.5,1,3,10,30},yticklabels={$0.5$,$1$,$3$,$10$,$30$}]
  \fill[band,opacity=0.72] (axis cs:-0.35,0.28) rectangle (axis cs:4.55,1.6);
  \node[anchor=south west,font=\footnotesize,inkPrim] at (axis cs:-0.28,0.42) {chem.\ accuracy};
  % REAL hardware recovery: MEAN over 8 recovery seeds on the SAME cached ibm_marrakesh counts
  % (job da125f2ein7c73bcsqs0, no new QPU); shaded envelope = +-1 s.d. across the 8 seeds at every step.
  \addplot[draw=none,name path=hwhi,forget plot] coordinates {(0,35.110) (1,6.958) (2,2.397) (3,1.199) (4,0.731)};
  \addplot[draw=none,name path=hwlo,forget plot] coordinates {(0,27.049) (1,5.267) (2,1.827) (3,0.877) (4,0.446)};
  \addplot[warmFill,fill opacity=0.30,forget plot] fill between[of=hwhi and hwlo];
  \addplot[simSlate,line width=1.4pt,dash pattern=on 5pt off 3pt,mark=square,mark size=2.6pt,
           mark options={draw=simSlate,fill=white,line width=1pt}] coordinates {(0,31.174) (1,29.600) (2,29.546) (3,29.546)};
  % mean over 8 seeds + a capped +-1 s.d. error bar at EVERY step, on top of the shaded envelope
  \addplot[hwWarm,line width=2.2pt,mark=*,mark size=3.0pt,mark options={fill=hwWarm,draw=white,line width=0.7pt},
     error bars/.cd,y dir=both,y explicit,error bar style={hwWarm,line width=1.1pt},
     error mark=-,error mark options={rotate=90,hwWarm,line width=1.1pt,mark size=4pt}]
     coordinates {(0,31.079)+-(0,4.030) (1,6.113)+-(0,0.845) (2,2.112)+-(0,0.285) (3,1.038)+-(0,0.161) (4,0.588)+-(0,0.142)};
  % which curve is which: a leader from each label lands ON its own curve; plus the two headline numbers
  \node[anchor=south,font=\footnotesize,inkPrim] (nl) at (axis cs:1.85,44)
    {noiseless (same circuit), $29.5$~mHa};
  \draw[-{Stealth[length=1.7mm]},inkPrim,line width=0.5pt] (nl.south) -- (axis cs:1.55,30.0);
  \node[anchor=south,align=center,font=\footnotesize,inkPrim] (nh) at (axis cs:3.20,5.2)
    {device counts\\$0.59\pm0.14$~mHa};
  \draw[-{Stealth[length=1.7mm]},inkPrim,line width=0.5pt] (nh.south) -- (axis cs:3.92,0.70);
\end{axis}
\end{tikzpicture}

\caption{\textbf{The complete quantum-hardware record of this work: two runs on IBM~Heron processors.}
\textbf{(a)}~Electron-addition branch of the local spectral function, $A^{+}_{0\uparrow}(\omega)$, of the
$L=6$ periodic Hubbard ring at $U/t=4$,
$\eta=0.15\,t$ ($12$ qubits; the classical validation of Fig.~\ref{fig:method} is a different instance,
$U/t=8$), computed on the subspace post-selected from $350{,}000$ computational-basis bitstrings---$50{,}000$ from each of seven
circuits---measured on \texttt{ibm\_fez} (measurement and Pauli
twirling and dynamical decoupling; no readout-error mitigation and no configuration recovery). The retained shots
covered the complete $(N{+}1)$ sector ($|\mathcal S|\!=\!300/300$), so the reconstruction is exact by
coverage; they were post-selected in reversed bit order, so every retained shot is a device error
(Sec.~\ref{sec:scope:hw}). The panel is an execution record, and neither a fidelity measurement nor a
claim of advantage.
\textbf{(b)}~Ground-state energy error of stretched $\mathrm{N_2}$ ($R=2.0$~\AA, CAS(10e,12o)/cc-pVDZ,
$24$ qubits) versus self-consistent configuration-recovery step, on \texttt{ibm\_marrakesh} ($10{,}000$ shots,
CCSD-seeded local unitary cluster-Jastrow (LUCJ) circuit). Recovery on the device counts (solid) reaches
$0.59\pm0.14$~mHa of the exact FCI value---well below chemical accuracy---whereas the noiseless statevector
simulation of the \emph{same} circuit (dashed) is stranded at $29.5$~mHa (its recorded recovery history
ends at step~$3$, unchanged since step~$2$). That contrast is not read here as device noise \emph{helping}: the
mechanism such a reading needs is decided by the two recovered subspace sizes, and the record
stores neither (Sec.~\ref{sec:scope:hw}). The
solid curve is the mean over $8$ recovery seeds on the cached device counts; the shaded envelope and the
capped whiskers at every step both show the $\pm1$ population s.d.\ (normalized by $8$, not $7$) across those $8$ seeds (widest at
step~$0$, $31.1\pm4.0$~mHa, where it starts together with the noiseless curve, and tapering as recovery
converges to $0.59\pm0.14$~mHa; best seed $0.40$~mHa). No random-bitstring control was run, so the
panel cannot separate what the device counts contributed from what recovery alone supplies; the
controlled recovery-versus-no-recovery comparison, under a simulated bit-flip model, is
Fig.~\ref{fig:noiserec}. These two runs are the entirety of the quantum
hardware used here; every other result in this work is an exact classical simulation.}\label{fig:heron}\end{figure*}
